% The Haskell 98 Foreign Function Interface
% [An Addendum to the Definition of Haskell 98]
%
% Editor: Manuel M T Chakravarty
%
% Copyright [2002..2003] Manuel M T Chakravarty
%
% The authors intend this Report to belong to the entire Haskell community, and
% so we grant permission to copy and distribute it for any purpose, provided
% that it is reproduced in its entirety, including this Notice.  Modified
% versions of this Report may also be copied and distributed for any purpose,
% provided that the modified version is clearly presented as such, and that it
% does not claim to be a definition of the Haskell 98 Foreign Function
% Interface. 

% Changes since RC15:
% * 6.3: Footnote regarding __STDC_ISO_10646__ added to text introducing
%        `CWString'.
%
% Changes since RC14:
% * 6.2: CWChar -> CWchar
% * 6.3: - CWChar -> CWchar
%        - Stated explicitly that memory allocated by `newCString' and friends
%          can be deallocated by `Foreign.Marshal.Alloc.free'
%        - Improved documentation
%
% Changes since RC13:
% * 5.3: Fixed typo
% * 5.7: Fixed a mistake in the type of `peekByteOff' and `pokeByteOff' (the
%        type variable constrained by `Storable' must be different from the
%        parameter of the `Ptr')
% * 6.3: Improved documentation
%
% Changes since RC12:
% * Acks : Added John Meacham
% * 4.1.5: Bug fix courtesy of Wolfgang Thaller
% * 5.5  : Added `FinalizerEnvPtr', `newForeignPtrEnv', and
%          `addForeignPtrFinalizerEnv'
% * 6.3  : Added John Meacham proposal for `wchar_t' support as well localised
%          string marshalling; in particular, this adds `CWString' and
%          `CWStringLen' as well as the `CWString' and the `CAString' family
%          of marshalling routines.  In addition, `charIsRepresentable' was
%          added. 
%
% Changes since RC11:
% * 5.5: Swapped argument order of `newForeignPtr' and `addForeignPtrFinalizer'
%
% Changes since RC10:
% * 3.3  : Clarified use of foreign functions of pure type
% * 4.1.1: Clarified the meaning of foreign imports without a \mbox{$\it &\it $} that have a
%          non-functional type in Haskell
% * 5.1  : Clarified the scope of safe use of unsafePerformIO
% * 5.5  : \mbox{$\it pre-emptive$} dropped in footnote regarding finalizers; added
%          `newForeignPointer_' and renamed `foreignPtrToPtr' to
%          `unsafeForeignPtrToPtr'
% * Typos throughout
%
% Changes since RC9:
% * 1:     Mentioning interaction with foreign threads as an open problem.
% * 2 & 3: Removed `threadsafe' again, as the proposal for thread support is
%          still evolving and it is not yet clear whether a new safety level
%          is required.
% * 5.5:   Added the type synonym `FinalizerPtr' and rewrote the documentation
%          of finalizers.
% * 5.6:   Clarified the description of `StablePtr'
% * 5.8:   Added `finalizerFree'
% * 6.2:   All the types in CTypes must be newtypes that are exported
%          abstractly. 
%
% Changes since RC8:
% * 5.8: `MarshallAlloc.reallocBytes' is no longer permitted on memory
%        allocated with `alloca' or `allocaBytes'. 
% * 6.1: Deinitialisation of the RTS via `hs_exit()' followed by
%        (re)initialisation with `hs_init()' must be supported.
%
% Changes since RC7:
% * Clarified the lexis of C identifiers and C header file names
% * In `ForeignPtr', added `mallocForeignPtrArray' and `mallocForeignPtrArray0'
% * Clarified spec of allocations functions adding constraints taken from the
%   corresponding C routines
% * `mallocBytes' and `allocaBytes' must align memory sufficiently for any
%   basic foreign type that fits into the allocated block
% * Removed typos in the description of the module `ForeignPtr'
% * Added Peter Gammie to the list of acknowledged people
% * `addForeignPtrFinalizer' guarantees that finalizers for a single foreign
%   pointer are executed in the opposite order as they were added.
% * `Storable': Require that the size is divisible by the alignment
% * Added Ross Paterson to the list of acknowledged people
% * Added hs_free_fun_ptr() and hs_free_stable_ptr()
% * Changed order of arguments of `mkIOError' and `annotateIOError' to match
%   with the current implementation in GHC's FFI libraries.
%
% Changes since RC6:
% * Fixed typos
%
% Changes since RC5:
% * Author list: changed Alastair Reid's institution
% * 1.4:   Clarified the wording
% * 4.1:   Explicitly stated that access to pre-processor symbols is not
%          provided by the FFI
% * 4.1.1: Removed [lib] from impent syntax and discussion
% * 4.1.3: Added parentheses round FunPtr ft to make it easier to 
%          understand a tolerably complex type.
% * 4.1.4: Removed all mention of library objects; clarified that header files
%          do not impact the semantics of foreign calls, but may be required
%          for correct code generation by some systems
% * 5.2:   Clarified that all operations in Bits are member functions of the
%          type class.  Reverse the meaning of the sign of the second argument
%          for `rotate' and `shift' (this makes it the same as GHC used all
%          the time).  `bitSize' on `Integer' etc is now undefined.
% * 5.5:   Finalisers must be external functions to facilitate the
%          implementation on Haskell systems that do not support pre-emptive
%          concurrency.
%          Added mallocForeignPtr and mallocForeignPtrBytes.
% * 6:     Specified that HsBool==int in table2
%          Relabelled column 1 in table 3 (C symbol -> CPP symbol)
%          Replaced 0 and 1 with HS_BOOL_FALSE/TRUE
% * 6.1:   Clarified that nullPtr (nullFunPtr) coincides with (HsPtr) NULL and
%          (HsFunPtr) NULL, respectively.
%          Allowing multiple calls to hs_init() and clarified the constraints
%          on the relative timing between hs_set_argv() and
%          getProgName/getArgs. 
%          Added hs_perform_gc().
%
% Changes since RC4:
% * 5.6: Clarified documentation of `StablePtr's (RC 5)
%
% Changes between RC2 and RC4:
%
% * 5.8: Clarified documentation for `MarshalAlloc.free'.
% * 5.8: Added `MarshalAlloc.realloc'.
% * 3: Added the new safety level `threadsafe' with an explanation at the end
%   of 3.3.
% * 3: Replaced the nontermional `entity' by `impent' and `expent' to
%   distinguish between import and export entities (as they are defined
%   differently in later sections).
% * 3.2: Clarified the description of foreign types; so far, `IO ()' was
%   strictly speaking not included as a valid return type.  Currently,
%   functions of type `a -> ()' are included.  Do we want this?  Their use
%   might not be portable if they include side effects.
% * 4.1.5: New section discussing the traps & pitfalls of type promotion with
%   C bindings.

% TODO:
% * Implement HTMLization.  (Malcolm suggests using
%   <http://pauillac.inria.fr/~maranget/hevea/>)

% TODO after Version 1.0:
%
% * Review suggestions by Antony Courtney <antony@apocalypse.org> re FFI
%   support for Java.

\chapter{Foreign Function Interface}
\label{ffi}

% Old FFI addendum text:
% The definition of Haskell 98~\cite{haskell98}, while being comprehensive with
% respect to the functional core language, does lack a range of features of more
% operational flavour, such as a foreign language interface, concurrency
% support, and fully fledged exception handling.  As these features are of
% central importance to many real world applications of the language, there is a
% danger that different implementations become de facto incompatible for such
% applications due to system-specific extensions of the core language.  The
% present FFI specification is aimed at reducing this risk by defining a simple,
% yet comprehensive extension to Haskell 98 for the purpose of interfacing to
% program components implemented in a language other than Haskell.

The Foreign Function Interface (FFI) has two purposes: it enables (1) to
describe in Haskell the interface to foreign language functionality and
(2) to use from foreign code Haskell routines.  More generally, its aim
is to support the implementation of programs in a mixture of Haskell and
other languages such that the source code is portable across different
implementations of Haskell and non-Haskell systems as well as
independent of the architecture and operating system.

% Old FFI addendum text:
% The design as presented in this report builds on experiences with a number of
% foreign function interfaces that, over time, have been provided by the major
% Haskell implementations.  Central in the final design was the goal to be
% comprehensive while being simple and minimising changes with respect to
% Haskell 98; the latter includes to avoid pollution of the name space with new
% keywords.  Consequently, as much as possible of the FFI functionality is
% realised in the form of libraries.  Simplicity generally overruled maximum
% convenience for the programmer as a design goal.  Thus, support for more
% convenient interface specifications is the domain of system-independent tools
% that generate code following the present specification.

\section{Foreign Languages}

The Haskell FFI currently only specifies the interaction between Haskell
code and foreign code that follows the C calling convention.  However,
the design of the FFI is such that it enables the modular extension of
the present definition to include the calling conventions of other
programming languages, such as C++ and Java.  A precise definition of
the support for those languages is expected to be included in later
versions of the language.  The second major omission is the definition
of the interaction with multithreading in the foreign language and, in
particular, the treatment of thread-local state, and so these details
are currently implementation-defined.

The core of the present specification is independent of the foreign language
that is used in conjunction with Haskell.  However, there are two areas where
FFI specifications must become language specific: (1) the specification of
external names and (2) the marshalling of the basic types of a foreign
language.  As an example of the former, consider that in C~\cite{C} a simple
identifier is sufficient to identify an object, while
Java~\cite{gosling-etal:Java}, in general, requires a qualified name in
conjunction with argument and result types to resolve possible overloading.
Regarding the second point, consider that many languages do not specify the
exact representation of some basic types.  For example the type \mbox{\tt int} in
C may be 16, 32, or 64 bits wide.  Similarly, Haskell guarantees
only that \mbox{\tt Int} covers at least the range \([-2^{29}, 2^{29} - 1]\) (Section~\ref{numbers}).  As
a consequence, to reliably represent values of C's \mbox{\tt int} in Haskell, we
have to introduce a new type \mbox{\tt CInt}, which is guaranteed to match the
representation of \mbox{\tt int}.

The specification of external names, dependent on a calling convention, is
described in Section~\ref{sec:extent}, whereas the marshalling of the basic
types in dependence on a foreign language is described in
Section~\ref{sec:marshalling}.

\section{Contexts}

For a given Haskell system, we define the \emph{Haskell context} to be the
execution context of the abstract machine on which the Haskell system is
based.  This includes the heap, stacks, and the registers of the abstract
machine and their mapping onto a concrete architecture.  We call any other
execution context an \emph{external context.}  Generally, we cannot assume any
compatibility between the data formats and calling conventions between the
Haskell context and a given external context, except where Haskell explicitly prescribes a specific data format.

The principal goal of a foreign function interface is to provide a
programmable interface between the Haskell context and external
contexts.  As a result Haskell threads can access data in external
contexts and invoke functions that are executed in an external context
as well as vice versa.  In the rest of this definition, external
contexts are usually identified by a calling convention.

\subsection{Cross Language Type Consistency}

Given that many external languages support static types, the question arises
whether the consistency of Haskell types with the types of the external
language can be enforced for foreign functions.  Unfortunately, this is, in
general, not possible without a significant investment on the part of the
implementor of the Haskell system (i.e., without implementing a dedicated type
checker).  For example, in the case of the C calling convention, the only
other approach would be to generate a C prototype from the Haskell type and
leave it to the C compiler to match this prototype with the prototype that is
specified in a C header file for the imported function.  However, the Haskell
type is lacking some information that would be required to pursue this route.
In particular, the Haskell type does not contain any information as to when
\mbox{\tt const} modifiers have to be emitted.  

As a consequence, this definition does not require the Haskell system to check
consistency with foreign types.  Nevertheless, Haskell systems are encouraged
to provide any cross language consistency checks that can be implemented with
reasonable effort.

\section{Lexical Structure}

The FFI reserves a single keyword \mbox{\tt foreign}, and a set of special
identifiers.  The latter have a special meaning only within foreign
declarations, but may be used as ordinary identifiers elsewhere.

The special identifiers \mbox{\tt ccall}, \mbox{\tt cplusplus}, \mbox{\tt dotnet},
\mbox{\tt jvm}, and \mbox{\tt stdcall} are defined to denote calling conventions.
However, a concrete implementation of the FFI is free to support additional,
system-specific calling conventions whose name is not explicitly listed here.

To refer to objects of an external C context, we introduce the following
phrases:
%
\begin{flushleft}\it\begin{tabbing}
\hspace{0.6in}\=\hspace{3.1in}\=\kill
$\it chname$\>\makebox[3.5em]{$\rightarrow$}$\it \{chchar\}\ \makebox{\tt .}\ \makebox{\tt h}$\>\makebox[3em]{}$\it (\tr{C\ header\ filename})$\\ 
$\it cid$\>\makebox[3.5em]{$\rightarrow$}$\it letter\ \{letter\ |\ ascDigit\}$\>\makebox[3em]{}$\it (\tr{C\ identifier})$\\ 
$\it chchar$\>\makebox[3.5em]{$\rightarrow$}$\it letter\ |\ ascSymbol_{\langle{}\makebox{\tt {\char'46}}\rangle{}}$\\ 
$\it letter$\>\makebox[3.5em]{$\rightarrow$}$\it ascSmall\ |\ ascLarge\ |\ \makebox{\tt {\char'137}}$
\end{tabbing}\end{flushleft}
%
The range of lexemes that are admissible for \mbox{$\it chname$} is a subset of
those permitted as arguments to the \mbox{\tt {\char'43}include} directive in C.  In
particular, a file name \mbox{$\it chname$} must end in the suffix \mbox{\tt .h}.  The
lexemes produced by \mbox{$\it cid$} coincide with those allowed as C identifiers,
as specified in~\cite{C}.

\section{Foreign Declarations}

The syntax of foreign declarations is as follows:
%
\begin{flushleft}\it\begin{tabbing}
\hspace{0.6in}\=\hspace{3.1in}\=\kill
$\it topdecl$\>\makebox[3.5em]{$\rightarrow$}$\it \makebox{\tt foreign}\ fdecl$\\ 
$\it fdecl$\>\makebox[3.5em]{$\rightarrow$}$\it \makebox{\tt import}\ callconv\ [safety]\ impent\ var\ \makebox{\tt ::}\ ftype$\>\makebox[3em]{}$\it (\tr{define\ variable})$\\ 
$\it $\>\makebox[3.5em]{$|$}$\it \makebox{\tt export}\ callconv\ expent\ var\ \makebox{\tt ::}\ ftype$\>\makebox[3em]{}$\it (\tr{expose\ variable})$\\ 
$\it callconv$\>\makebox[3.5em]{$\rightarrow$}$\it \makebox{\tt ccall}\ |\ \makebox{\tt stdcall}\ |\ \makebox{\tt cplusplus}$\>\makebox[3em]{}$\it (\tr{calling\ convention})$\\ 
$\it $\>\makebox[3.5em]{$|$}$\it \makebox{\tt jvm}\ |\ \makebox{\tt dotnet}$\\ 
$\it $\>\makebox[3.5em]{$|$}$\it \mbox{\bf\ system-specific\ calling\ conventions}$\\ 
$\it impent$\>\makebox[3.5em]{$\rightarrow$}$\it [string]$\\ 
$\it expent$\>\makebox[3.5em]{$\rightarrow$}$\it [string]$\\ 
$\it safety$\>\makebox[3.5em]{$\rightarrow$}$\it \makebox{\tt unsafe}\ |\ \makebox{\tt safe}$
\end{tabbing}\end{flushleft}
%
There are two flavours of foreign declarations: import and export
declarations.  An import declaration makes an \emph{external entity,} i.e., a
function or memory location defined in an external context, available in the
Haskell context.  Conversely, an export declaration defines a function of the
Haskell context as an external entity in an external context.  Consequently,
the two types of declarations differ in that an import declaration defines a
new variable, whereas an export declaration uses a variable that is already
defined in the Haskell module.

The external context that contains the external entity is determined by the
calling convention given in the foreign declaration.  Consequently, the exact
form of the specification of the external entity is dependent on both the
calling convention and on whether it appears in an import declaration (as
\mbox{$\it impent$}) or in an export declaration (as \mbox{$\it expent$}).  To provide
syntactic uniformity in the presence of different calling conventions, it is
guaranteed that the description of an external entity lexically appears as a
Haskell string lexeme.  The only exception is where this string would be the
empty string (i.e., be of the form \mbox{\tt ""}); in this case, the string may be
omitted in its entirety.

\subsection{Calling Conventions}
\label{sec:call-conv}

The binary interface to an external entity on a given architecture is
determined by a calling convention.  It often depends on the programming
language in which the external entity is implemented, but usually is more
dependent on the system for which the external entity has been compiled.

As an example of how the calling convention is dominated by the system rather
than the programming language, consider that an entity compiled to byte code
for the Java Virtual Machine (JVM)~\cite{lindholm-etal:JVM} needs to be
invoked by the rules of the JVM rather than that of the source language in
which it is implemented (the entity might be implemented in Oberon, for
example).

Any implementation of the Haskell FFI must at least implement the C calling
convention denoted by \mbox{\tt ccall}.  All other calling conventions are
optional.  Generally, the set of calling conventions is open, i.e., individual
implementations may elect to support additional calling conventions.  In
addition to \mbox{\tt ccall}, Table~\ref{tab:callconv} specifies a range of
identifiers for common calling conventions.
%
\begin{table}[tbp]
  \begin{center}
    \begin{tabular}{|l|l|}
      \hline
      Identifier & Represented calling convention\\
      \hline\hline
      \mbox{\tt ccall} 
      & Calling convention of the standard C compiler on a system\\
      \mbox{\tt cplusplus}
      & Calling convention of the standard C{+}{+} compiler on a system\\
      \mbox{\tt dotnet}
      & Calling convention of the \textsc{.net} platform\\
      \mbox{\tt jvm} 
      & Calling convention of the Java Virtual Machine\\
      \mbox{\tt stdcall}
      & Calling convention of the Win32 API (matches Pascal conventions)\\
      \hline
    \end{tabular}
    \caption{Calling conventions}
    \label{tab:callconv}
  \end{center}
\end{table}
%
Implementations need not implement all of these conventions, but if any is
implemented, it must use the listed name.  For any other calling convention,
implementations are free to choose a suitable name.

Only the semantics of the calling conventions \mbox{\tt ccall} and \mbox{\tt stdcall} are
defined herein; more calling conventions may be added in future versions
of Haskell.

It should be noted that the code generated by a Haskell system to implement a
particular calling convention may vary widely with the target code of that
system.  For example, the calling convention \mbox{\tt jvm} will be trivial to
implement for a Haskell compiler generating Java code, whereas for a Haskell
compiler generating C code, the Java Native Interface (JNI)~\cite{liang:JNI}
has to be targeted.

\subsection{Foreign Types}
\label{sec:foreign-types}

The following types constitute the set of \emph{basic foreign types}:
%
\begin{itemize}
\item \mbox{\tt Char}, \mbox{\tt Int}, \mbox{\tt Double}, \mbox{\tt Float}, and \mbox{\tt Bool} as
  exported by the Haskell \mbox{\tt Prelude} as well as
\item \mbox{\tt Int8}, \mbox{\tt Int16}, \mbox{\tt Int32}, \mbox{\tt Int64}, \mbox{\tt Word8},
  \mbox{\tt Word16}, \mbox{\tt Word32}, \mbox{\tt Word64}, \mbox{\tt Ptr\ a}, \mbox{\tt FunPtr\ a},
  and \mbox{\tt StablePtr\ a}, for any type \mbox{\tt a}, as exported by \mbox{\tt Foreign}
  (Section~\ref{module:Foreign}).
\end{itemize}
%
A Haskell system that implements the FFI needs to be able to pass these types
between the Haskell and the external context as function arguments and
results.

Foreign types are produced according to the following grammar:
%
\begin{flushleft}\it\begin{tabbing}
\hspace{0.6in}\=\hspace{3.1in}\=\kill
$\it ftype$\>\makebox[3.5em]{$\rightarrow$}$\it frtype$\\ 
$\it $\>\makebox[3.5em]{$|$}$\it fatype\ \makebox{\tt ->}\ ftype$\\ 
$\it frtype$\>\makebox[3.5em]{$\rightarrow$}$\it fatype$\\ 
$\it $\>\makebox[3.5em]{$|$}$\it \makebox{\tt ()}$\\ 
$\it fatype$\>\makebox[3.5em]{$\rightarrow$}$\it qtycon\ atype_1\ \ldots\ atype_k$\>\makebox[3em]{}$\it (k\ \geq\ 0)$
\end{tabbing}\end{flushleft}
%
A foreign type is the Haskell type of an external entity.  Only a subset of
Haskell's types are permissible as foreign types, as only a restricted set of
types can be canonically transferred between the Haskell context and an
external context.  A foreign type has the form
\[
\textit{at}_1\mbox{\tt \ ->\ }\cdots\mbox{\tt \ ->\ }\textit{at}_n\mbox{\tt \ ->\ }\textit{rt}
\]
where \(n\geq0\).  It implies that the arity of the external entity is $n$.

External functions are strict in all arguments.

\paragraph{Marshallable foreign types.}

The argument types \(\textit{at}_i\) produced by \mbox{$\it fatype$} must be
\emph{marshallable foreign types;} that is, either

\begin{itemize}
\item a basic foreign type,

\item a type synonym that expands to a marshallable foreign type,

\item a type \mbox{$\it T\ t'_1\ \ldots \ t'_n$} where \mbox{$\it T$} is defined by a \mbox{\tt newtype} declaration

\begin{quote}
\mbox{\tt newtype\ }\mbox{$\it T\ a_1\ \ldots \ a_n$}\mbox{\tt \ =\ }\mbox{$\it N\ t$}
\end{quote}

and

\begin{itemize}
\item the constructor \mbox{$\it N$} is visible where \mbox{$\it T$} is used,

\item \mbox{$\it t[t'_1/a_1..t'_n/a_n]$} is a marshallable foreign type
\end{itemize}
\end{itemize}

Consequently, in order for a type defined by \mbox{\tt newtype} to be used in a
\mbox{\tt foreign} declaration outside of the module that defines it, the type
must not be exported abstractly.  The module \mbox{\tt Foreign.C.Types} that
defines the Haskell equivalents for C types follows this convention;
see Chapter~\ref{module:Foreign.C.Types}.

\paragraph{Marshallable foreign result types.}

The result type \textit{rt} produced by \mbox{$\it frtype$} must be a
\emph{marshallable foreign result type;} that is, either

\begin{itemize}
\item the type \mbox{\tt ()},
\item a type matching \mbox{\tt Prelude.IO\ }$t$, where $t$ is a marshallable foreign type or \mbox{\tt ()},
\item a basic foreign type,
\item a type synonym that expands to marshallable foreign result type,
\item a type \mbox{$\it T\ t'_1\ \ldots \ t'_n$} where \mbox{$\it T$} is defined by a \mbox{\tt newtype} declaration

\begin{quote}
\mbox{\tt newtype\ }\mbox{$\it T\ a_1\ \ldots \ a_n$}\mbox{\tt \ =\ }\mbox{$\it N\ t$}
\end{quote}

and

\begin{itemize}
\item the constructor \mbox{$\it N$} is visible where \mbox{$\it T$} is used,

\item \mbox{$\it t[t'_1/a_1..t'_n/a_n]$} is a marshallable foreign result type
\end{itemize}
\end{itemize}

\subsection{Import Declarations}
\label{sec:import}

Generally, an import declaration has the form
%
\[
\mbox{\tt foreign}~\mbox{\tt import}~c~e~v~\mbox{\tt ::}~t
\]
%
which declares the variable $v$ of type $t$ to be defined externally.
Moreover, it specifies that $v$ is evaluated by executing the external entity
identified by the string $e$ using calling convention $c$.  The precise form
of $e$ depends on the calling convention and is detailed in
Section~\ref{sec:extent}.  If a variable $v$ is defined by an import
declaration, no other top-level declaration for $v$ is allowed in the same
module.  For example, the declaration
\bprog
{\small\begin{verbatim}
foreign import ccall "string.h strlen"
   cstrlen :: Ptr CChar -> IO CSize
\end{verbatim}}
\eprog
introduces the function \mbox{\tt cstrlen}, which invokes the external function
\mbox{\tt strlen} using the standard C calling convention.  Some external entities
can be imported as pure functions; for example,
\bprog
{\small\begin{verbatim}
foreign import ccall "math.h sin"
   sin :: CDouble -> CDouble.
\end{verbatim}}
\eprog
%
Such a declaration asserts that the external entity is a true function; i.e.,
when applied to the same argument values, it always produces the same result.

Whether a particular form of external entity places a constraint on the
Haskell type with which it can be imported is defined in
Section~\ref{sec:extent}.  Although, some forms of external entities restrict
the set of Haskell types that are permissible, the system can generally not
guarantee the consistency between the Haskell type given in an import
declaration and the argument and result types of the external entity.  It is
the responsibility of the programmer to ensure this consistency.

Optionally, an import declaration can specify, after the calling convention,
the safety level that should be used when invoking an external entity.  A
\mbox{\tt safe} call is less efficient, but guarantees to leave the Haskell system
in a state that allows callbacks from the external code.  In contrast, an
\mbox{\tt unsafe} call, while carrying less overhead, must not trigger a callback
into the Haskell system.  If it does, the system behaviour is undefined.  The
default for an invocation is to be \mbox{\tt safe}.  Note that a callback into
the Haskell system implies that a garbage collection might be triggered after
an external entity was called, but before this call returns.  Consequently,
objects other than stable pointers (cf.\ Section~\ref{module:Foreign.StablePtr}) may be
moved or garbage collected by the storage manager.

\subsection{Export Declarations}

The general form of export declarations is
%
\[
\mbox{\tt foreign}~\mbox{\tt export}~c~e~v~\mbox{\tt ::}~t
\]
%
Such a declaration enables external access to $v$, which may be a value, field
name, or class method that is declared on the top-level of the same module or
imported.  Moreover, the Haskell system defines the external entity described
by the string $e$, which may be used by external code using the calling
convention $c$; an external invocation of the external entity $e$ is
translated into evaluation of $v$.  The type $t$ must be an instance of the
type of $v$.  For example, we may have
\bprog
{\small\begin{verbatim}
foreign export ccall "addInt"   (+) :: Int   -> Int   -> Int
foreign export ccall "addFloat" (+) :: Float -> Float -> Float
\end{verbatim}}
\eprog
If an evaluation triggered by an external invocation of an exported Haskell
value returns with an exception, the system behaviour is undefined.  Thus,
Haskell exceptions have to be caught within Haskell and explicitly marshalled
to the foreign code.


\section{Specification of External Entities}
\label{sec:extent}

Each foreign declaration has to specify the external entity that is accessed
or provided by that declaration.  The syntax and semantics of the notation
that is required to uniquely determine an external entity depends heavily on
the calling convention by which this entity is accessed.  For example, for the
calling convention \mbox{\tt ccall}, a global label is sufficient.  However, to
uniquely identify a method in the calling convention \mbox{\tt jvm}, type
information has to be provided.  For the latter, there is a choice between the
Java source-level syntax of types and the syntax expected by JNI---but,
clearly, the syntax of the specification of an external entity depends on the
calling convention and may be non-trivial.

Consequently, the FFI does not fix a general syntax for denoting external
entities, but requires both \mbox{$\it impent$} and \mbox{$\it expent$} to take the
form of a Haskell \mbox{$\it string$} literal.  The formation rules for the values
of these strings depend on the calling convention and a Haskell system
implementing a particular calling convention will have to parse these strings
in accordance with the calling convention.

Defining \mbox{$\it impent$} and \mbox{$\it expent$} to take the form of a
\mbox{$\it string$} implies that all information that is needed to statically
analyse the Haskell program is separated from the information needed to
generate the code interacting with the foreign language.  This is, in
particular, helpful for tools processing Haskell source code.  When ignoring
the entity information provided by \mbox{$\it impent$} or \mbox{$\it expent$}, foreign
import and export declarations are still sufficient to infer identifier
definition and use information as well as type information.

For more complex calling conventions, there is a choice between the user-level
syntax for identifying entities (e.g., Java or C{+}{+}) and the system-level
syntax (e.g., the type syntax of JNI or mangled C{+}{+}, respectively).  If
such a choice exists, the user-level syntax is preferred.  Not only because it
is more user friendly, but also because the system-level syntax may not be
entirely independent of the particular implementation of the foreign language.

The following defines the syntax for specifying external entities and their
semantics for the calling conventions \mbox{\tt ccall} and \mbox{\tt stdcall}.  Other
calling conventions from Table~\ref{tab:callconv} are expected to be defined
in future versions of Haskell.


\subsection{Standard C Calls}
\label{sec:ccall}

The following defines the structure of external entities for foreign
declarations under the \mbox{\tt ccall} calling convention for both import and
export declarations separately.  Afterwards additional constraints on the type
of foreign functions are defined.

The FFI covers only access to C functions and global variables.  There are no
mechanisms to access other entities of C programs.  In particular, there is no
support for accessing pre-processor symbols from Haskell, which includes
\mbox{\tt {\char'43}define}d constants.  Access from Haskell to such entities is the
domain of language-specific tools, which provide added convenience over the
plain FFI as defined here.

\paragraph{Import Declarations}

For import declarations, the syntax for the specification of external entities
under the \mbox{\tt ccall} calling convention is as follows:
\begin{flushleft}\it\begin{tabbing}
\hspace{0.6in}\=\hspace{3.1in}\=\kill
$\it impent$\>\makebox[3.5em]{$\rightarrow$}$\it \makebox{\tt "}\ [\makebox{\tt static}]\ [chname]\ [\makebox{\tt {\char'46}}]\ [cid]\ \makebox{\tt "}$\>\makebox[3em]{}$\it (\tr{static\ function\ or\ address})$\\ 
$\it $\>\makebox[3.5em]{$|$}$\it \makebox{\tt "}\ dynamic\ \makebox{\tt "}$\>\makebox[3em]{}$\it (\tr{stub\ factory\ importing\ addresses})$\\ 
$\it $\>\makebox[3.5em]{$|$}$\it \makebox{\tt "}\ wrapper\ \makebox{\tt "}$\>\makebox[3em]{}$\it (\tr{stub\ factory\ exporting\ thunks})$
\end{tabbing}\end{flushleft}
The first alternative either imports a static function \mbox{$\it cid$} or, if
\mbox{\tt {\char'46}} precedes the identifier, a static address.  If \mbox{$\it cid$} is
omitted, it defaults to the name of the imported Haskell variable.  The
optional filename \mbox{$\it chname$} specifies a C header file, where the
intended meaning is that the header file declares the C entity identified by
\mbox{$\it cid$}.  In particular, when the Haskell system compiles Haskell to C
code, the directive
%
\begin{quote}
  \mbox{\tt {\char'43}include\ "}\mbox{$\it chname$}\mbox{\tt "}
\end{quote}
%
needs to be placed into any generated C file that refers to the foreign entity
before the first occurrence of that entity in the generated C file.

The second and third alternative, identified by the keywords \mbox{\tt dynamic}
and \mbox{\tt wrapper}, respectively, import stub functions that have to be
generated by the Haskell system.  In the case of \mbox{\tt dynamic}, the stub
converts C function pointers into Haskell functions; and conversely, in the
case of \mbox{\tt wrapper}, the stub converts Haskell thunks to C function
pointers.  If neither of the specifiers \mbox{\tt static}, \mbox{\tt dynamic}, or
\mbox{\tt wrapper} is given, \mbox{\tt static} is assumed.  The specifier
\mbox{\tt static} is nevertheless needed to import C routines that are named
\mbox{\tt dynamic} or \mbox{\tt wrapper}.

It should be noted that a static foreign declaration that does not import an
address (i.e., where \mbox{\tt {\char'46}} is not used in the specification of the external
entity) always refers to a C function, even if the Haskell type is
non-functional.  For example, 
\bprog
{\small\begin{verbatim}
foreign import ccall foo :: CInt
\end{verbatim}}
\eprog
refers to a pure C function \mbox{\tt foo} with no arguments that returns an
integer value.  Similarly, if the type is \mbox{\tt IO\ CInt}, the declaration
refers to an impure nullary function.  If a Haskell program needs to access a
C variable \mbox{\tt bar} of integer type,
\bprog
{\small\begin{verbatim}
foreign import ccall "&" bar :: Ptr CInt
\end{verbatim}}
\eprog
must be used to obtain a pointer referring to the variable.  The variable can
be read and updated using the routines provided by the module \mbox{\tt Foreign.Storable}
(cf.\ Section~\ref{module:Foreign.Storable}).

\paragraph{Export Declarations}

External entities in \mbox{$\it ccall$} export declarations are of the form
%
\begin{flushleft}\it\begin{tabbing}
\hspace{0.6in}\=\hspace{3.1in}\=\kill
$\it expent$\>\makebox[3.5em]{$\rightarrow$}$\it \makebox{\tt "}\ [cid]\ \makebox{\tt "}$
\end{tabbing}\end{flushleft}
%
The optional C identifier \mbox{$\it cid$} defines the external name by which the
exported Haskell variable is accessible in C.  If it is omitted, the external
name defaults to the name of the exported Haskell variable.

\paragraph{Constraints on Foreign Function Types}

In the case of import declaration, there are, depending on the kind of import
declaration, constraints regarding the admissible Haskell type that the
variable defined in the import may have.  These constraints are specified in
the following.
%
\begin{description}
\item[Static Functions.]  A static function can be of any foreign type; in
  particular, the result type may or may not be in the IO monad.  If a
  function that is not pure is not imported in the IO monad, the system
  behaviour is undefined.  Generally, no check for consistency with the C type
  of the imported label is performed.

  As an example, consider
  %
\bprog
{\small\begin{verbatim}
foreign import ccall "static stdlib.h"
   system :: Ptr CChar -> IO CInt
\end{verbatim}}
\eprog
  %
  This declaration imports the \mbox{\tt system()} function whose prototype is
  available from \mbox{\tt stdlib.h}.

\item[Static addresses.]  The type of an imported address is constrained to be
  of the form \mbox{\tt Ptr\ }\textit{a} or \mbox{\tt FunPtr\ }\textit{a}, where
  \textit{a} can be any type.

  As an example, consider
  %
\bprog
{\small\begin{verbatim}
foreign import ccall "errno.h &errno" errno :: Ptr CInt
\end{verbatim}}
\eprog
  %
  It imports the address of the variable \mbox{\tt errno}, which is of the C type
  \mbox{\tt int}.

\item[Dynamic import.]  The type of a \mbox{$\it dynamic$} stub has to be of the
  form \mbox{\tt (FunPtr\ }\textit{ft}\mbox{\tt )\ ->\ }\textit{ft}, where \textit{ft} may
  be any foreign type.

  As an example, consider
  %
\bprog
{\small\begin{verbatim}
foreign import ccall "dynamic" 
  mkFun :: FunPtr (CInt -> IO ()) -> (CInt -> IO ())
\end{verbatim}}
\eprog
  %
  The stub factory \mbox{\tt mkFun} converts any pointer to a C function that gets
  an integer value as its only argument and does not have a return value into
  a corresponding Haskell function.

\item[Dynamic wrapper.]  The type of a \mbox{$\it wrapper$} stub has to be of the
  form \textit{ft}\mbox{\tt \ ->\ IO\ (FunPtr\ }\textit{ft}), where
  \textit{ft} may be any foreign type.

  As an example, consider
  %
\bprog
{\small\begin{verbatim}
foreign import ccall "wrapper" 
  mkCallback :: IO () -> IO (FunPtr (IO ()))
\end{verbatim}}
\eprog
  %
  The stub factory \mbox{\tt mkCallback} turns any Haskell computation of type
  \mbox{\tt IO\ ()} into a C function pointer that can be passed to C routines,
  which can call back into the Haskell context by invoking the referenced
  function.

\end{description}

\paragraph{Specification of Header Files}

A C header specified in an import declaration is always included by
\mbox{\tt {\char'43}include\ "}\mbox{$\it chname$}\mbox{\tt "}.  There is no explicit support for
\mbox{\tt {\char'43}include\ <}\mbox{$\it chname$}\mbox{\tt >} style inclusion.  The ISO C99~\cite{C99}
standard guarantees that any search path that would be used for a
\mbox{\tt {\char'43}include\ <}\mbox{$\it chname$}\mbox{\tt >} is also used for \mbox{\tt {\char'43}include\ "}\mbox{$\it chname$}\mbox{\tt "} and it is guaranteed that these paths are searched after
all paths that are unique to \mbox{\tt {\char'43}include\ "}\mbox{$\it chname$}\mbox{\tt "}.  Furthermore,
we require that \mbox{$\it chname$} ends in \mbox{\tt .h} to make parsing of the
specification of external entities unambiguous.
  
The specification of include files has been kept to a minimum on purpose.
Libraries often require a multitude of include directives, some of which may
be system-dependent.  Any design that attempts to cover all possible
configurations would introduce significant complexity.  Moreover, in the
current design, a custom include file can be specified that uses the standard
C preprocessor features to include all relevant headers.

Header files have no impact on the semantics of a foreign call, and whether an
implementation uses the header file or not is implementation-defined.
However, as some implementations may require a header file that supplies a
correct prototype for external functions in order to generate correct code,
portable FFI code must include suitable header files.

\paragraph{C Argument Promotion}

The argument passing conventions of C are dependent on whether a function
prototype for the called functions is in scope at a call site.  In particular,
if no function prototype is in scope, \emph{default argument promotion} is
applied to integral and floating types.  In general, it cannot be expected
from a Haskell system that it is aware of whether a given C function was
compiled with or without a function prototype being in scope.  For the sake of
portability, we thus require that a Haskell system generally implements calls
to C functions as well as C stubs for Haskell functions as if a function
prototype for the called function is in scope.

This convention implies that the onus for ensuring the match between C and
Haskell code is placed on the FFI user.  In particular, when a C function that
was compiled without a prototype is called from Haskell, the Haskell signature
at the corresponding \mbox{\tt foreign\ import} declaration must use the types
\emph{after} argument promotion.  For example, consider the following C
function definition, which lacks a prototype:
%
\bprog
{\small\begin{verbatim}
void foo (a)
float a;
{
  ...
}
\end{verbatim}}
\eprog
%
The lack of a prototype implies that a C compiler will apply default argument
promotion to the parameter \mbox{\tt a}, and thus, \mbox{\tt foo} will expect to
receive a value of type \mbox{\tt double}, \emph{not} \mbox{\tt float}.  Hence, the
correct \mbox{\tt foreign\ import} declaration is
%
\bprog
{\small\begin{verbatim}
foreign import ccall foo :: Double -> IO ()
\end{verbatim}}
\eprog

In contrast, a C function compiled with the prototype
%
\bprog
{\small\begin{verbatim}
void foo (float a);
\end{verbatim}}
\eprog
%
requires
%
\bprog
{\small\begin{verbatim}
foreign import ccall foo :: Float -> IO ()
\end{verbatim}}
\eprog

A similar situation arises in the case of \mbox{\tt foreign\ export} declarations
that use types that would be altered under the C default argument promotion
rules.  When calling such Haskell functions from C, a function prototype
matching the signature provided in the \mbox{\tt foreign\ export} declaration must
be in scope; otherwise, the C compiler will erroneously apply the promotion
rules to all function arguments.

Note that for a C function defined to accept a variable number of arguments,
all arguments beyond the explicitly typed arguments suffer argument promotion.
However, because C permits the calling convention to be different for such
functions, a Haskell system will, in general, not be able to make use of
variable argument functions.  Hence, their use is deprecated in portable code.


\subsection{Win32 API Calls}

The specification of external entities under the \mbox{\tt stdcall} calling
convention is identical to that for standard C calls.  The two calling
conventions only differ in the generated code.


% \begin{FUTURE} % ===== Material for future extension =======================
% 
% \subsection{C{+}{+} Calls}
% 
% The syntax for the specification of external entities under the
% \mbox{\tt cplusplus} calling convention is
% 
% \subsection{JVM Calls}
% 
% The syntax for the specification of external entities under the \mbox{\tt jvm}
% calling convention is 
% %
% 
%   \grule{impent}{%
%     "jtype\mbox{$\it $}jqid\mbox{$\it ($}jtypes\mbox{$\it )$}}
%   \gor[constructor call]{%
%     \mbox{$\it new$}jqid\mbox{$\it ($}jtypes\mbox{$\it )$}}
%   \grule[$n\geq0$]{jtypes}{%
%     \gnterm[1]{jtype},\gellipse,\gnterm[n]{jtype}}
% 
% %
% where \mbox{$\it jqid$} is a qualified Java identifier and \mbox{$\it jtype$} a Java
% types as defined in~\cite{gosling-etal:Java}.
% 
% \begin{verbatim}
% FIXME: 
% - force the inclusion of the return type in case of \mbox{$\it new$}?
% \end{verbatim}
% 
% \subsection{.NET Calls}
% 
% The syntax for the specification of external entities under the \mbox{\tt dotnet}
% calling convention is
% 
% \end{FUTURE}% =============================================================


\section{Marshalling}
\label{sec:marshalling}

In addition to the language extension discussed in previous sections, the FFI
includes a set of standard libraries, which ease portable use of foreign
functions as well as marshalling of compound structures.  Generally, the
marshalling of Haskell structures into a foreign representation and vice versa
can be implemented in either Haskell or the foreign language.  At least where
the foreign language is at a significantly lower level, e.g.\ C, there are
good reasons for doing the marshalling in Haskell:
%
\begin{itemize}
\item Haskell's lazy evaluation strategy would require any foreign code that
  attempts to access Haskell structures to force the evaluation of these
  structures before accessing them. This would lead to complicated code in the
  foreign language, but does not need any extra consideration when coding the
  marshalling in Haskell.
\item Despite the fact that marshalling code in Haskell tends to look like C
  in Haskell syntax, the strong type system still catches many errors that
  would otherwise lead to difficult-to-debug runtime faults.
\item Direct access to Haskell heap structures from a language like
  C---especially, when marshalling from C to Haskell, i.e., when Haskell
  structures are created---carries the risk of corrupting the heap, which
  usually leads to faults that are very hard to debug.
\end{itemize}
%
Consequently, the Haskell FFI emphasises Haskell-side marshalling.

The interface to the marshalling libraries is provided by the module
\mbox{\tt Foreign} (Chapter~\ref{module:Foreign}) plus a language-dependent module per supported language.  In
particular, the standard requires the availability of the module
\mbox{\tt Foreign.C} (Chapter~\ref{module:Foreign.C}), which simplifies portable interfacing with external C code.
Language-dependent modules, such as \mbox{\tt Foreign.C}, generally provide Haskell
types representing the basic types of the foreign language using a
representation that is compatible with the foreign types as implemented by the
default implementation of the foreign language on the present architecture.
This is especially important for languages where the standard leaves some
aspects of the implementation of basic types open.  For example, in C, the
size of the various integral types is not fixed.  Thus, to represent C
interfaces faithfully in Haskell, for each integral type in C, we need to have
an integral type in Haskell that is guaranteed to have the same size as the
corresponding C type.

\section{The External C Interface}

\begin{table}
  \begin{center}
    \begin{tabular}{|l|l|l|}
      \hline
      C symbol          & Haskell symbol & Constraint on concrete C type\\
      \hline\hline
      \mbox{\tt HsChar}     & \mbox{\tt Char}    
      & integral type\\
      \hline
      \mbox{\tt HsInt}      & \mbox{\tt Int}
      & signed integral type, $\geq30$ bit\\
      \hline
      \mbox{\tt HsInt8}     & \mbox{\tt Int8}
      & signed integral type, 8 bit; \mbox{\tt int8{\char'137}t} if available\\
      \hline
      \mbox{\tt HsInt16}    & \mbox{\tt Int16}
      & signed integral type, 16 bit; \mbox{\tt int16{\char'137}t} if available\\
      \hline
      \mbox{\tt HsInt32}    & \mbox{\tt Int32}
      & signed integral type, 32 bit; \mbox{\tt int32{\char'137}t} if available\\
      \hline
      \mbox{\tt HsInt64}    & \mbox{\tt Int64}
      & signed integral type, 64 bit; \mbox{\tt int64{\char'137}t} if available\\ 
      \hline
      \mbox{\tt HsWord8}    & \mbox{\tt Word8}
      & unsigned integral type, 8 bit; \mbox{\tt uint8{\char'137}t} if available\\
      \hline
      \mbox{\tt HsWord16}   & \mbox{\tt Word16}
      & unsigned integral type, 16 bit; \mbox{\tt uint16{\char'137}t} if available\\
      \hline
      \mbox{\tt HsWord32}   & \mbox{\tt Word32}
      & unsigned integral type, 32 bit; \mbox{\tt uint32{\char'137}t} if available\\
      \hline
      \mbox{\tt HsWord64}   & \mbox{\tt Word64}
      & unsigned integral type, 64 bit; \mbox{\tt uint64{\char'137}t} if available\\
      \hline
      \mbox{\tt HsFloat}    & \mbox{\tt Float}
      & floating point type\\
     \hline
      \mbox{\tt HsDouble}   & \mbox{\tt Double}
      & floating point type\\
     \hline
      \mbox{\tt HsBool}     & \mbox{\tt Bool}
      & \mbox{\tt int}\\
     \hline
      \mbox{\tt HsPtr}      & \mbox{\tt Ptr\ a}
      & \mbox{\tt (void\ *)}\\
     \hline
      \mbox{\tt HsFunPtr}   & \mbox{\tt FunPtr\ a}
      & \mbox{\tt (void\ (*)(void))}\\
     \hline
      \mbox{\tt HsStablePtr}& \mbox{\tt StablePtr\ a}
      & \mbox{\tt (void\ *)}\\
     \hline
    \end{tabular}
    \caption{C Interface to Basic Haskell Types}
    \label{tab:c-haskell-types}
  \end{center}
\end{table}
%
\begin{table}
  \begin{center}
%    \begin{tabular}{|l|l|l|}
    \begin{tabular}{|l|l|p{30ex}|}
      \hline
      CPP symbol           & Haskell value & Description\\
      \hline\hline
      \mbox{\tt HS{\char'137}CHAR{\char'137}MIN} & \mbox{\tt minBound\ ::\ Char}
      & \\
      \hline
      \mbox{\tt HS{\char'137}CHAR{\char'137}MAX} & \mbox{\tt maxBound\ ::\ Char}
      & \\
      \hline
      \mbox{\tt HS{\char'137}INT{\char'137}MIN} & \mbox{\tt minBound\ ::\ Int}
      & \\
      \hline
      \mbox{\tt HS{\char'137}INT{\char'137}MAX} & \mbox{\tt maxBound\ ::\ Int}
      & \\
      \hline
      \mbox{\tt HS{\char'137}INT8{\char'137}MIN} & \mbox{\tt minBound\ ::\ Int8}
      & \\
      \hline
      \mbox{\tt HS{\char'137}INT8{\char'137}MAX} & \mbox{\tt maxBound\ ::\ Int8}
      & \\
      \hline
      \mbox{\tt HS{\char'137}INT16{\char'137}MIN} & \mbox{\tt minBound\ ::\ Int16}
      & \\
      \hline
      \mbox{\tt HS{\char'137}INT16{\char'137}MAX} & \mbox{\tt maxBound\ ::\ Int16}
      & \\
      \hline
      \mbox{\tt HS{\char'137}INT32{\char'137}MIN} & \mbox{\tt minBound\ ::\ Int32}
      & \\
      \hline
      \mbox{\tt HS{\char'137}INT32{\char'137}MAX} & \mbox{\tt maxBound\ ::\ Int32}
      & \\
      \hline
      \mbox{\tt HS{\char'137}INT64{\char'137}MIN} & \mbox{\tt minBound\ ::\ Int64}
      & \\
      \hline
      \mbox{\tt HS{\char'137}INT64{\char'137}MAX} & \mbox{\tt maxBound\ ::\ Int64}
      & \\
      \hline
      \mbox{\tt HS{\char'137}WORD8{\char'137}MAX} & \mbox{\tt maxBound\ ::\ Word8}
      & \\
      \hline
      \mbox{\tt HS{\char'137}WORD16{\char'137}MAX} & \mbox{\tt maxBound\ ::\ Word16}
      & \\
      \hline
      \mbox{\tt HS{\char'137}WORD32{\char'137}MAX} & \mbox{\tt maxBound\ ::\ Word32}
      & \\
      \hline
      \mbox{\tt HS{\char'137}WORD64{\char'137}MAX} & \mbox{\tt maxBound\ ::\ Word64}
      & \\
      \hline
      \mbox{\tt HS{\char'137}FLOAT{\char'137}RADIX} & \mbox{\tt floatRadix\ ::\ Float}
      & \\
      \hline
      \mbox{\tt HS{\char'137}FLOAT{\char'137}ROUND} & n/a
      & rounding style as per~\cite{C99}\\
      \hline
      \mbox{\tt HS{\char'137}FLOAT{\char'137}EPSILON} & n/a
      & difference between 1 and the least value greater
      than 1 as per~\cite{C99}\\
      \hline
      \mbox{\tt HS{\char'137}DOUBLE{\char'137}EPSILON} & n/a
      & (as above)\\
      \hline
      \mbox{\tt HS{\char'137}FLOAT{\char'137}DIG} & n/a
      & number of decimal digits as per~\cite{C99}\\
      \hline
      \mbox{\tt HS{\char'137}DOUBLE{\char'137}DIG} & n/a
      & (as above)\\
      \hline
      \mbox{\tt HS{\char'137}FLOAT{\char'137}MANT{\char'137}DIG} & \mbox{\tt floatDigits\ ::\ Float}
      & \\
      \hline
      \mbox{\tt HS{\char'137}DOUBLE{\char'137}MANT{\char'137}DIG} & \mbox{\tt floatDigits\ ::\ Double}
      & \\
      \hline
      \mbox{\tt HS{\char'137}FLOAT{\char'137}MIN} & n/a
      & minimum floating point number as per~\cite{C99}\\
      \hline
      \mbox{\tt HS{\char'137}DOUBLE{\char'137}MIN} & n/a
      & (as above)\\
      \hline
      \mbox{\tt HS{\char'137}FLOAT{\char'137}MIN{\char'137}EXP} & \mbox{\tt fst\ .\ floatRange\ ::\ Float}
      & \\
      \hline
      \mbox{\tt HS{\char'137}DOUBLE{\char'137}MIN{\char'137}EXP} & \mbox{\tt fst\ .\ floatRange\ ::\ Double}
      & \\
      \hline
      \mbox{\tt HS{\char'137}FLOAT{\char'137}MIN{\char'137}10{\char'137}EXP} & n/a
      & minimum decimal exponent as per~\cite{C99}\\
      \hline
      \mbox{\tt HS{\char'137}DOUBLE{\char'137}MIN{\char'137}10{\char'137}EXP} & n/a
      & (as above)\\
      \hline
      \mbox{\tt HS{\char'137}FLOAT{\char'137}MAX} & n/a
      & maximum floating point number as per~\cite{C99}\\
      \hline
      \mbox{\tt HS{\char'137}DOUBLE{\char'137}MAX} & n/a
      & (as above)\\
      \hline
      \mbox{\tt HS{\char'137}FLOAT{\char'137}MAX{\char'137}EXP} & \mbox{\tt snd\ .\ floatRange\ ::\ Float}
      & \\
      \hline
      \mbox{\tt HS{\char'137}DOUBLE{\char'137}MAX{\char'137}EXP} & \mbox{\tt snd\ .\ floatRange\ ::\ Double}
      & \\
      \hline
      \mbox{\tt HS{\char'137}FLOAT{\char'137}MAX{\char'137}10{\char'137}EXP} & n/a
      & maximum decimal exponent as per~\cite{C99}\\
      \hline
      \mbox{\tt HS{\char'137}DOUBLE{\char'137}MAX{\char'137}10{\char'137}EXP} & n/a
      & (as above)\\
      \hline
      \mbox{\tt HS{\char'137}BOOL{\char'137}FALSE} & False
      & \\
      \hline
      \mbox{\tt HS{\char'137}BOOL{\char'137}TRUE} & True
      & \\
      \hline
    \end{tabular}
    \caption{C Interface to Range and Precision of Basic Types}
    \label{tab:c-haskell-values}
  \end{center}
\end{table}
%
Every Haskell system that implements the FFI needs to provide a C header file
named \mbox{\tt HsFFI.h} that defines the C symbols listed in
Tables~\ref{tab:c-haskell-types} and~\ref{tab:c-haskell-values}.
Table~\ref{tab:c-haskell-types} table lists symbols that represent types
together with the Haskell type that they represent and any constraints that
are placed on the concrete C types that implement these symbols.  When a C
type \mbox{\tt HsT} represents a Haskell type \mbox{\tt T}, the occurrence of \mbox{\tt T}
in a foreign function declaration should be matched by \mbox{\tt HsT} in the
corresponding C function prototype.  Indeed, where the Haskell system
translates Haskell to C code that invokes \mbox{\tt foreign} \mbox{\tt import}ed C
routines, such prototypes need to be provided and included via the header that
can be specified in external entity strings for foreign C functions (cf.\ 
Section~\ref{sec:ccall}); otherwise, the system behaviour is undefined.  It is
guaranteed that the Haskell value \mbox{\tt nullPtr} is mapped to \mbox{\tt (HsPtr)\ NULL} in C and \mbox{\tt nullFunPtr} is mapped to \mbox{\tt (HsFunPtr)\ NULL} and
vice versa.

Table~\ref{tab:c-haskell-values} contains symbols characterising the range and
precision of the types from Table~\ref{tab:c-haskell-types}.  Where available,
the table states the corresponding Haskell values.  All C symbols, with the
exception of \mbox{\tt HS{\char'137}FLOAT{\char'137}ROUND} are constants that are suitable for use in
\mbox{\tt {\char'43}if} preprocessing directives.  Note that there is only one rounding
style (\mbox{\tt HS{\char'137}FLOAT{\char'137}ROUND}) and one radix (\mbox{\tt HS{\char'137}FLOAT{\char'137}RADIX}), as
this is all that is supported by ISO C~\cite{C99}.

Moreover, an implementation that does not support 64 bit integral types on the
C side should implement \mbox{\tt HsInt64} and \mbox{\tt HsWord64} as a structure.  In
this case, the bounds \mbox{\tt HS{\char'137}INT64{\char'137}MIN}, \mbox{\tt HS{\char'137}INT64{\char'137}MAX}, and
\mbox{\tt HS{\char'137}WORD64{\char'137}MAX} are undefined.

In addition, to the symbols from Table~\ref{tab:c-haskell-types}
and~\ref{tab:c-haskell-values}, the header \mbox{\tt HsFFI.h} must also contain
the following prototypes:
%
\begin{quote}
\begin{verbatim}
void hs_init     (int *argc, char **argv[]);
void hs_exit     (void);
void hs_set_argv (int argc, char *argv[]);

void hs_perform_gc (void);

void hs_free_stable_ptr (HsStablePtr sp);
void hs_free_fun_ptr    (HsFunPtr fp);
\end{verbatim}
\end{quote}
%
These routines are useful for mixed language programs, where the main
application is implemented in a foreign language that accesses routines
implemented in Haskell.  The function \mbox{\tt hs{\char'137}init()} initialises the
Haskell system and provides it with the available command line arguments.
Upon return, the arguments solely intended for the Haskell runtime system are
removed (i.e., the values that \mbox{\tt argc} and \mbox{\tt argv} point to may have
changed).  This function must be called during program startup before any
Haskell function is invoked; otherwise, the system behaviour is undefined.
Conversely, the Haskell system is deinitialised by a call to
\mbox{\tt hs{\char'137}exit()}.  Multiple invocations of \mbox{\tt hs{\char'137}init()} are permitted,
provided that they are followed by an equal number of calls to
\mbox{\tt hs{\char'137}exit()} and that the first call to \mbox{\tt hs{\char'137}exit()} is after the
last call to \mbox{\tt hs{\char'137}init()}.  In addition to nested calls to
\mbox{\tt hs{\char'137}init()}, the Haskell system may be de-initialised with
\mbox{\tt hs{\char'137}exit()} and be re-initialised with \mbox{\tt hs{\char'137}init()} at a later
point in time.  This ensures that repeated initialisation due to multiple
libraries being implemented in Haskell is covered.

The Haskell system will ignore the command line arguments passed to the second
and any following calls to \mbox{\tt hs{\char'137}init()}.  Moreover, \mbox{\tt hs{\char'137}init()} may
be called with \mbox{\tt NULL} for both \mbox{\tt argc} and \mbox{\tt argv}, signalling
the absence of command line arguments.

The function \mbox{\tt hs{\char'137}set{\char'137}argv()} sets the values returned by the functions
\mbox{\tt getProgName} and \mbox{\tt getArgs} of the module \mbox{\tt System.Environment} (Section~\ref{module:System.Environment}).  This function may only be invoked after
\mbox{\tt hs{\char'137}init()}.  Moreover, if \mbox{\tt hs{\char'137}set{\char'137}argv()} is called at all, this
call must precede the first invocation of \mbox{\tt getProgName} and
\mbox{\tt getArgs}.  Note that the separation of \mbox{\tt hs{\char'137}init()} and
\mbox{\tt hs{\char'137}set{\char'137}argv()} is essential in cases where in addition to the Haskell
system other libraries that process command line arguments during
initialisation are used.

The function \mbox{\tt hs{\char'137}perform{\char'137}gc()} advises the Haskell storage manager to
perform a garbage collection, where the storage manager makes an effort to
releases all unreachable objects.  This function must not be invoked from C
functions that are imported \mbox{\tt unsafe} into Haskell code nor may it be used
from a finalizer.

Finally, \mbox{\tt hs{\char'137}free{\char'137}stable{\char'137}ptr()} and \mbox{\tt hs{\char'137}free{\char'137}fun{\char'137}ptr()} are
the C counterparts of the Haskell functions \mbox{\tt freeStablePtr} and
\mbox{\tt freeHaskellFunPtr}.
